%%%%%%%%%%%%%%%%%%%%%%%%%%%%%%%%%%%%%%%%%%%%%%%%%%%
% Signal Processing Laboratory (LTS5) - EPFL      %
% LaTeX beamposter template                       %
% Authors:                                        %
%   D. Perdios – dimitris.perdios@epfl.ch         %
%   A. Besson – adrien.besson@epfl.ch             %
% v0.1 - 08.05.17                                 %
% Typeset configuration: pdfLaTeX + Biber         %
%%%%%%%%%%%%%%%%%%%%%%%%%%%%%%%%%%%%%%%%%%%%%%%%%%%


% New command used to create a blank column

\newcommand{\AddBlankColumn}[1]{
	\begin{column}{#1}\end{column}
}

%%% SIZES OF THE MARGINS AND COLUMNS

%%% CUSTOM MAKETITLE COMMAND %%%
\newlength{\LogoWidthLeft}
\settowidth{\LogoWidthLeft}{\includegraphics[height=\LogoHeight]{logos/lts5_logo.pdf}}
\newlength{\LogoWidthRight}
\settowidth{\LogoWidthRight}{\includegraphics[height=\LogoHeight]{logos/epfl_logo.pdf}}
\setlength{\titlewidth}{\LeftColumnWidth + \RightColumnWidth + \BlankColumnWidthCenter- \LogoWidthLeft - \LogoWidthRight - \BlankColumnWidthLeft-\BlankColumnWidthRight}
\newsavebox{\titlebox}
\newlength{\myposterheight}
\renewcommand{\maketitle}{%
	\savebox{\titlebox}{%
		\hspace{\BlankColumnWidthLeft}
		\begin{minipage}{\LogoWidthLeft}%
			\includegraphics[height=\LogoHeight]{logos/basp_logo.png}
			\vspace{\LogoSkip}
			\centering
			\includegraphics[height=\LogoHeight]{logos/lts5_logo.pdf}
		\end{minipage} 
		\begin{minipage}{\titlewidth}%
			\centering
			\LARGE\inserttitle \leavevmode\\%
			\normalsize\insertauthor \leavevmode\\%
			\small\insertinstitute %
		\end{minipage} 
		\begin{minipage}{\LogoWidthRight}%
			\includegraphics[height=\LogoHeight]{logos/epfl_logo.pdf}
			\vspace{\LogoSkip}
			\centering
			\includegraphics[height=\LogoHeight]{logos/lts5_logo.pdf}
		\end{minipage}
		\hspace{\BlankColumnWidthRight}
	}%
	\settoheight{\titleheight}{\usebox{\titlebox}}
	\setlength{\myposterheight}{0.97\textheight - \titleheight - \TitleColumnSpace}
	\usebox{\titlebox}
}

% *** USEFUL COMMANDS ***
\newcommand{\ie}{\textit{i.e.}}
\newcommand{\etal}{\textit{et al.}}
\newcommand{\vect}[1]{\bm{#1}}
\newcommand{\ser}[2]{#1^#2}
\newcommand{\mat}[1]{\mathsf{#1}}
\DeclareMathOperator{\R}{\mathbb{R}}